% ============================================================
%  SECTION 1: INTRODUCTION
% ============================================================

\section{Introduction}

\textcite{dynarski2000hope} showed that Georgia's HOPE Scholarship
widened the gap in college attendance between higher- and lower-income
students. The finding crystallized a distributional critique of merit
aid: because eligibility is conditioned on academic achievement, which
correlates with family income, broad-based merit programs
disproportionately subsidize students who would attend college regardless
\parencite{dynarski2004merit, jones2022merit}. This view now dominates
the policy debate.

A systematic review of 71 quasi-experimental studies finds that
merit-based grants rarely improve outcomes for disadvantaged students and
can increase inequality \parencite{herbaut2020works}. We document an
exception.

Using tax-linked income data from \textcite{chetty2020mobility}, we
estimate how Tennessee's HOPE Scholarship reshaped the full parental
income distribution at public colleges. The results are the opposite of
what the conventional view predicts: following HOPE implementation,
public colleges experienced a \emph{progressive} compositional shift.
Bottom-quintile students gained share at public institutions, while
top-quintile students lost share. The effect is monotonic across all
five income quintiles---Q1 through Q3 gain, Q4 and Q5 lose---and every
coefficient is statistically significant. This monotonic gradient is the
paper's central empirical result.

Aggregate attendance effects and institutional composition effects are
distinct estimands that need not move in the same direction. In a
two-sector model---attend college or not---individual enrollment
responses fully determine institutional composition. But when students
choose among public, private, and out-of-state alternatives, the link
breaks: the \emph{aggregate} attendance gap can widen (because
higher-income students are more likely to qualify) while the
\emph{sector-specific} composition of public colleges simultaneously
improves (because the students who actually sort into the public sector
are drawn disproportionately from the bottom of the income
distribution).\footnote{This is a form of Simpson's paradox: the
aggregate relationship between merit aid and income composition reverses
sign when conditioned on enrollment sector. The aggregate effect reflects
who attends college; the composition effect reflects who attends
\emph{which} college.} The reconciliation depends on where displaced
students go. If top-quintile students exit to out-of-state or private
alternatives rather than entering public colleges, the public sector
becomes more progressive even as overall college-going becomes more
unequal. We formalize this insight in a three-destination sorting model
(public, private, out-of-state) and test its predictions directly.

Institutional composition matters independently of individual enrollment
decisions. The peer effects literature demonstrates that student body
composition shapes academic outcomes
\parencite{sacerdote2001peer}. \textcite{chetty2020mobility} show that
access to high-mobility institutions drives intergenerational economic
outcomes, establishing that who attends which institution has first-order
welfare consequences. If merit aid systematically alters the income
profile of public college cohorts, it reshapes the peer environment and
mobility prospects for all enrolled students---not only those on the
margin of attendance.

Four pieces of evidence are consistent with a sorting interpretation.
First, total enrollment at Tennessee public colleges is unchanged after
HOPE. The share shifts reflect real reallocation---11 additional
bottom-quintile students per cohort enter, 34 top-quintile students per
cohort exit---not a denominator artifact. Second, Tennessee private
colleges show no mirror image of the public-sector shift: the decline in
top-quintile representation at public colleges is not offset by gains at
Tennessee privates, consistent with out-of-state exit as the relevant
margin. Third, two-year and four-year institutions adjust through
different channels. Community colleges exhibit displacement
(top-quintile students replaced by bottom-quintile students, total seats
unchanged), while four-year colleges show expansion (enrollment grows,
composition shifts). Fourth, within-state comparisons restricted to
Tennessee institutions confirm that the composition change is specific
to the public sector.

These patterns may not be idiosyncratic to Tennessee. Our theoretical
framework proposes that program design and institutional context jointly
shape the direction of compositional sorting, through three parameters:
the eligibility threshold relative to the income distribution, the size
of the public-private subsidy differential, and interactions with
existing need-based aid. Georgia's
HOPE, which initially offset Pell Grant awards dollar-for-dollar and
imposed an income cap, eliminated the access margin for the
lowest-income families and muted the public-private subsidy
differential---conditions under which the model predicts regressive
sorting. Dynarski's findings for Georgia are consistent with this
prediction, as are the enrollment patterns documented by
\textcite{cornwell2006enrollment}. Tennessee's HOPE, which has no Pell
offset, no income cap, and a large public-private subsidy
differential, satisfies the conditions for progressive sorting. The same
framework explains both outcomes.

We test these predictions out of sample. Massachusetts' Adams
Scholarship \parencite{cohodes2014merit}, which shares Tennessee's
design features (moderate threshold, public-sector subsidy
concentration), independently produces a progressive shift of similar
magnitude ($+1.38$ pp for bottom-quintile share). South Dakota's
Opportunity Scholarship, which imposes a high ACT threshold and small
award, produces a strongly regressive shift---as the model predicts.
West Virginia's PROMISE, despite similar design parameters to
Tennessee, shows no compositional effect, highlighting that program
design is necessary but not sufficient: enrollment contraction at WV
public colleges overwhelmed any sorting dynamics.

We estimate treatment effects using \textcite{callaway2021difference} as
the primary estimator, with TWFE, \textcite{sun2021estimating}, and
\textcite{dechaisemartin2020two} for robustness. A bordering-state control
group passes pre-trend tests that the national comparison fails at
longer horizons. Results are stable across three treatment timing
definitions and two alternative control group specifications.
Under the relative-magnitude restriction of
\textcite{rambachan2023more}, the estimated effect survives violations
up to twice the magnitude observed in the pre-treatment period. Under
the stricter smoothness restriction, results are more fragile. We
interpret the evidence as credible but not definitive, and report all
sensitivity analyses transparently.

Our work connects three literatures. First, a large body of research
documents the enrollment effects of merit aid
\parencite{dynarski2000hope, dynarski2004merit, bruce2014jackpot,
cohodes2014merit, song2024hope, cummings2022race,
verzyl2023institutional}. This literature
focuses on individual enrollment decisions, not institutional
composition. Second, studies of policy-induced sorting show that
admissions and financing reforms can substantially reshape who attends
which institutions \parencite{bleemer2022affirmative,
card2005affirmative, hinrichs2012affirmative}. Third,
\textcite{chetty2020mobility} demonstrate that access to
high-mobility institutions shapes intergenerational economic outcomes,
establishing that institutional composition has first-order welfare
consequences. We bridge these literatures by documenting how merit aid is
associated with changes in the full income distribution across
institutional sectors.

This paper makes three contributions. First, we provide the first
evidence on how merit aid reshapes the \emph{full income distribution}
at public institutions, using tax-linked data that measures parental
income across all quintiles rather than relying on Pell eligibility as a
binary proxy. The monotonic quintile gradient---all five coefficients
significant and perfectly ordered---has not been documented in any prior
study of merit aid. Second, we develop a three-destination sorting
framework that reconciles the progressive institutional composition
effects we find with the regressive attendance effects documented by
\textcite{dynarski2000hope}, showing that the apparent contradiction
dissolves once student destinations are disaggregated. Third, we
propose that program design parameters---Pell interaction, subsidy
differential, threshold location---shape whether merit aid is associated
with progressive or regressive sorting, and test the framework across
four states: Tennessee and Massachusetts (progressive, as predicted),
South Dakota (regressive, as predicted), and West Virginia (null,
confounded by enrollment contraction).

The paper proceeds as follows. Section~\ref{sec:background} describes
the Tennessee HOPE Scholarship and the Georgia program that serves as
our theoretical comparison case. Section~\ref{sec:theory} presents the
three-destination sorting model and derives testable predictions.
Section~\ref{sec:data} describes the data.
Section~\ref{sec:empirical_strategy} discusses the empirical strategy.
Section~\ref{sec:results} presents results.
Section~\ref{sec:robustness} reports robustness checks.
Section~\ref{sec:discussion} discusses cross-state evidence from
Georgia, Massachusetts, South Dakota, and West Virginia.
Section~\ref{sec:conclusion} concludes.
